\begin{table}[H]
\centering
\caption{Estimated energy price elasticities, Spain}
\label{tab:es-epf-energy-elasticities}
\small
\setlength{\tabcolsep}{5pt}
\renewcommand{\arraystretch}{1.05}
\begin{tabular}{lcccc}
\toprule
Energy good & Price variable & Elasticity & Standard error & Observations \\
\midrule
Electricity & Household unit value & -1.005 & 0.011 & 36,769 \\
Electricity & Local price proxy & -0.718 & 0.049 & 36,769 \\
Gas & Household unit value & -1.076 & 0.021 & 20,014 \\
Gas & Local price proxy & -1.028 & 0.257 & 20,014 \\
Motor fuels & Household unit value & -1.777 & 0.063 & 19,805 \\
Motor fuels & Local price proxy & -0.701 & 0.510 & 19,805 \\
\bottomrule
\end{tabular}
\begin{minipage}{0.92\linewidth}
\footnotesize
\emph{Notes:} The dependent variable is log energy quantity. Estimates use
Spanish EPF microdata for available years 2020-2021. The household-unit-value
specifications use $\log(GASTO/CANTIDAD)$ as the price variable and include
autonomous-community and year fixed effects when available. The local-proxy
specifications use the median unit value by autonomous community, year, and
energy good; region fixed effects are not included because they would absorb
the identifying cross-regional price variation. All specifications control for
log total expenditure, log equivalised income, and household size. Source: INE
EPF microdata, authors' calculations.
\end{minipage}
\end{table}
